\subsection*{\theappendixitem.\ p-环的存在性}
\phantomsection
\addcontentsline{toc}{subsection}{\theappendixitem.\ p-环的存在性}

命题 4. --- 设 \(p\) 是一个素数，\(k\) 是一个特征为 \(p\) 的域，并设 \(n\) 是一个整数 \(\geqslant 1\) ，或为 \(+\infty\) 。则存在一个 \(p\) -环（§ 2，第 1 目，定义 1），其长度为 \(n\) ，剩余域同构于 \(k\) 。

可以把 \(k\) 看作 \(\mathbf{Z}/p\mathbf{Z}\) （局部环 \(\mathbf{Z}_{(p)}\) 的剩余域）的一个扩张。由定理 1 的推论，存在一个局部环 \(\mathbf{B}\) ，它是 \(\mathbf{Z}_{(p)}\) 的一个膨胀，使得 \(\kappa_{\mathrm{B}}\) 同构于 \(k\) 。局部环 \(\mathbf{Z}_{(p)}\) 是正则的，且 \(\{p\}\) 是 \(\mathbf{Z}_{(p)}\) 的一个坐标系。由第 2 目命题 2 的推论，环 \(\mathbf{B}\) 是正则的，且 \(\{p1_{\mathrm{B}}\}\) 是 \(\mathbf{B}\) 的一个坐标系。换言之，\(\mathbf{B}\) 是一个离散赋值环，其极大理想为 \(p\mathbf{B}\) 。于是完备化 \(\mathbf{C}\) （由 \(\mathbf{B}\) 完备化而得）是一个无限长度的 \(p\) -环，且剩余域 \(\kappa_{\mathrm{C}}\) 同构于 \(\kappa_{\mathrm{B}}\) ，从而同构于 \(k\) 。此外，对每个整数 \(n \geqslant 1\) ，C/ \(p^{n}\) C 是一个 \(p\) -环，其长度为 \(n\) ，剩余域同构于 \(\kappa_{\mathrm{C}}\) ，从而同构于 \(k\) 。

